\begin{unnumbered}

\section{Вступ}
    \textbf{Актуальність теми. } Сучасний етап розвитку професійного спорту  характеризується тотальною цифровізацією та
інтеграцією інтелектуальних  технологій у тренувальний та змагальний процеси. У волейболі, який є однією з
найдинамічніших командних ігор у світі, швидкість польоту м'яча при  атакуючих ударах може перевищувати 130 км/год,
а тривалість окремих ігрових фаз вимірюється частками секунди. За таких умов традиційний  візуальний аналіз з боку
тренерського штабу або використання простих систем двовимірної відеофіксації вже не задовольняють вимогам високоточного
розбору тактико-технічних дій.

    Автоматизований збір кінематичних метрик, таких як тривимірна траєкторія польоту м'яча, координати точок відскоку,
кути атаки та пікова швидкість польоту м'яча, дозволяє об'єктивно оцінювати ефективність дій гравців, моделювати
поведінку суперника та оптимізувати біомеханіку рухів під
час тренувань. Проте чинні комерційні рішення (наприклад, Hawk-Eye або Foxtenn) є монопольними системами з високою
вартістю впровадження, що вимагають встановлення від 10 до 20 синхронізованих високошвидкісних індустріальних камер
по всьому периметру залу. Це робить їх абсолютно недоступними для аматорських ліг, дитячо-юнацьких спортивних шкіл
та клубів середнього рівня.

    У зв'язку з цим виникає гостра потреба в розробці доступної, локальної та архітектурно гнучкої системи комп'ютерного
зору, здатної здійснювати точний тривимірний трекінг та аналітику ігрових параметрів волейболу на основі обмеженої
кількості джерел відеоданих (аж до однієї камери споживацького класу, наприклад, сучасного смартфона). Дана задача
лежить на перетині глибокого навчання (для детекції дрібних швидкісних об'єктів), проективної геометрії (для
калібрування простору) та теорії оптимальної фільтрації й оцінювання станів (для компенсації шумів вимірювання).
Реалізація такої системи в межах спеціальності "Комп'ютерна інженерія" дозволить створити оптимізований
програмно-алгоритмічний пайплайн, що ефективно поєднує низькорівневу обробку відеопотоків із сучасними методами
Data Science.

    \textbf{Мета дослідження. } Метою бакалаврської роботи є розробка та дослідження системи комп'ютерного зору
для високоточного автоматизованого аналізу волейбольних епізодів у тривимірному просторі за допомогою однієї
статичної камери на основі моделей глибокого навчання та нелінійного адаптивного фільтрування.

    Для досягнення поставленої мети необхідно вирішити такі завдання:
    \begin{itemize}
        \item проаналізувати існуючі методи детекції, трекінгу та відновлення
        3D-координат швидкісних об'єктів за монокулярними
        відеопотоками;
        \item обґрунтувати вибір та провести донавчання згорткової нейронної
        мережі сімейства YOLO для стабільної детекції волейбольного
        м'яча в умовах сильного розмиття рухом;
        \item розробити програмний модуль геометричного калібрування
        тривимірного простору майданчика на основі математичного
        апарату перспективного проектування камери(PnP);
        \item створити алгоритм зворотного проектування для
        побудови 3D-променів від оптичного центру лінзи через піксельні
        координати детекції та оцінки глибини польоту м'яча;
        \item дослідити проблему високочастотного тремтіння меж рамки
        детекції та спроектувати математичний апарат
        фільтрації траєкторії на основі нелінійного фільтра Калмана(\gls{UKF});
        \item імплементувати багатомодельну адаптивну архітектуру(IMM) для точного відстеження м'яча під час зміни фаз
        польоту, включаючи вільну балістичну параболу, момент ударного
        імпульсу та відскок від поверхонь;
        \item експериментально оцінити точність розрахунку миттєвої та пікової
        швидкостей польоту м'яча та провести статистичний аналіз
        тестових ігрових епізодів.
    \end{itemize}

    \textbf{Об'єкт дослідження} — процес відстеження просторового переміщення
волейбольного м'яча та визначення його кінематичних параметрів за допомогою
відеоданих у тривимірному просторі.
Предмет дослідження — методи комп'ютерного зору, моделі глибокого
навчання для об'єктної детекції, алгоритми проективної геометрії, математичні
методи нелінійної фільтрації станів динамічних систем та адаптивного
оцінювання.


    \textbf{Методи дослідження. }
Для вирішення поставлених завдань використано:
    \begin{itemize}
        \item методи глибокого навчання та штучних нейронних мереж (для виявлення м'яча
        на кадрах);
        \item апарат лінійної алгебри та проективної геометрії (для калібрування
        камери за реперними точками розмітки корту та реалізації Ray Casting);
        \item методи теорії ймовірностей та математичної статистики (для опису коваріаційних
        матриць шумів процесу та вимірювань);
        \item теорію оптимального керування та
        оцінювання динамічних станів, зокрема Unscented Kalman Filter (для
        згладжування траєкторій в умовах нелінійної аеродинаміки) та Interacting
        Multiple Models (для перемикання режимів руху).
    \end{itemize}

    \textbf{Наукова новизна. }
Наукова новизна отриманих результатів полягає у
розробці інтегрованого підходу до монокулярного 3D-трекінгу надшвидкісних
спортивних об'єктів, який, на відміну від існуючих методів оцінки глибини за
геометричними розмірами рамок детекції (що схильні до критичного
накопичення похибок), поєднує безпохідне нелінійне згладжування за
допомогою алгоритму UKF із динамічним перемиканням моделей руху (IMM).
Це дозволяє стабілізувати обчислення кінематичних параметрів у тривимірному
просторі без використання багатокамерних стереосистем.
Практичне значення. Розроблена система реалізована у вигляді
оптимізованого програмного комплексу мовою Python із використанням
сучасних індустріальних бібліотек (OpenCV, Ultralytics YOLO, FilterPy).
Спроектований пайплайн дозволяє проводити глибокий аналіз тренувальних та
змагальних відеозаписів, отриманих зі звичайних смартфонів, що забезпечує
низький поріг входження для спортивних клубів та індивідуальних тренерів.
Отримані аналітичні дані (графіки швидкостей, карти ударів, траєкторії) можуть
бути використані для автоматизації збору статистики та коригування техніки
виконання атакуючих елементів.
\end{unnumbered}

